% Imporditud: tools/import_eksperimendid.py
% Allikas: Eesti füüsikaolümpiaadi eksperimendi ülesannete
%   kogu (Rael Kalda, 2023), ülesanne Ü55, lahendus L55.
\setAuthor{EFO žürii}
\setRound{piirkonnavoor}
\setYear{2010}
\setNumber{G E1}
\setDifficulty{2}
\setTopic{staatika}
\setVahendid{tuntud massiga joonlaud (m = 10 g), tikutops koos tikkudega}
\setPunktid{10}
\setAllikas{Eksperimendi kogu Ü55, lahendus lk 57}
\setLahendusseis{toores}

\prob{Tiku mass}
Leidke ühe tiku mass.

\hint

\solu
Esmalt määrame joonlaua massikeskme asukoha. Selleks jagame skaala pooleks või tasakaalustame joonlaua tiku või laua serval ja paneme kirja tiku asukoha näidu joonlaual. (1 p) Järgmisena määrame tühja tikutopsi massi. Asetame tühja tikutopsi joonlauale võimalikult otsa lähedale ja leiame topsi keskkoha asukoha. Viimase saamiseks tuleks võtta topsi mõlema serva näidud ja arvutada aritmeetiline keskmine. (1 p) Tasakaalustame joonlaua ja paneme kirja toetuspunkti asukoha. (1 p) Saadud tulemuste põhjal arvutame joonlaua massikeskme ja tikutopsi massikeskme kaugused toetuspunktist $l_1$ ja $l_2$. Kangi seaduse tõttu:

m1gl1 = m2gl2 (2 p)

ja tühja toosi massiks saame

m$^2$ = m1l1 (1 p). $l_2$

Kordame katset täis tikutopsiga. (2 p) Ühe tiku massi leidmiseks tuleb täis tiku-

topsi massist lahutada tühja topsi mass ning jagada saadud tulemus topsis olnud

tikkude arvuga: mtais $-$ mtuhi n mtikk = (1 p).

Realistlik vastus: mtikk = 0,1 g (1 p).
\probend
